% LOI TUA --- Quyển XXV V4
\chapter*{Lời Tựa}
\addcontentsline{toc}{chapter}{Lời Tựa}
\markboth{Lời Tựa}{Lời Tựa}

\begin{truyen}{Lời Tựa}{The Classroom Multiverse}
\chuhoa{L}{ớp học} không chỉ là đường một chiều từ bục giảng xuống bàn học. Lớp học ngày nay là một ngã tư đường, nơi tư duy của người lớn va chạm với sự nhạy bén của Gen Z, nơi những kỳ vọng của cha mẹ đè nặng lên vai con trẻ, và nơi chính thầy cô cũng phải học cách thừa nhận những "lý sự cùn" của riêng mình. 
\textit{(A classroom is not a one-way street from the podium to the desks. Today's classroom is an intersection where adult thinking collides with Gen Z's sharpness, where parental expectations weigh on children, and where teachers themselves must learn to admit their own flawed logic.)}

Quyển V4 này mang tên \textbf{Đa Giác Học Đường}. Khác với những tập trước chỉ bóc mẽ học sinh, cuốn sách này chia thành 4 phần để nhìn nhận lại trường học từ nhiều góc độ. Chúng ta sẽ thấy những lý lẽ "thao túng tâm lý" của học trò, những lần giáo viên "dính trấu" vì luật lệ của chính mình, cuộc chiến ngầm với phụ huynh, và cả những ảo tưởng buồn cười về tương lai.
\textit{(This V4 is named \textbf{The Classroom Multiverse}. Unlike previous volumes that only exposed students, this book is divided into 4 parts to look at school from multiple angles. We will see the "mental warfare" logic of students, the times teachers get caught in their own rules, the silent wars with parents, and hilarious delusions about the future.)}

Mong rằng, V4 không chỉ mang lại tiếng cười châm biếm thường thấy, mà còn để lại những cái thở dài thấu hiểu. Bởi suy cho cùng, ai cũng có cái lý (cùn) của riêng mình!
\textit{(Hopefully, V4 will bring not only the usual satirical laughter but also sighs of empathy. Because in the end, everyone has their own (blunt) logic!)}
\end{truyen}

\begin{baihoc}
Trường học là một xã hội thu nhỏ. Bạn không thể phán xét một đứa trẻ nếu không nhìn vào áp lực từ cha mẹ chúng, và không thể trách một người thầy nếu không hiểu sự khác biệt thế hệ.
\textit{(School is a microcosm of society. You cannot judge a child without looking at the pressure from their parents, and you cannot blame a teacher without understanding the generational gap.)}
\end{baihoc}

\begin{ghinhoanh}
\textbf{Short English to remember:}

1. Empathy goes both ways. (Sự thấu cảm là con đường hai chiều).
2. The classroom is a multiverse of colliding expectations. (Lớp học là đa vũ trụ của những kỳ vọng va chạm nhau).
\end{ghinhoanh}

\begin{tuhocnhanh}
1. Bạn đã bao giờ dùng từ "chữa lành" để ngụy biện cho sự lười biếng chưa? \textit{(Have you ever used "healing" to excuse your laziness?)}
2. Phụ huynh bạn có bao giờ bênh vực bạn một cách vô lý trước mặt thầy cô? \textit{(Have your parents ever defended you unreasonably in front of teachers?)}
\end{tuhocnhanh}

\begin{flushright}
\textbf{Tôi Nguyễn Văn Sang}\\
\textit{GV THPT Nguyễn Hữu Cảnh - TP HCM}
\end{flushright}

\ngancach
